\chapter{Introduction}

This chapter introduces the research topic, provides an overview and motivation for the study, and outlines the main objectives and expected outcomes. It also outlines the methodology and briefly describes how the rest of the report is organized.

\section{Project Overview}
Ordinary Differential Equations (ODEs) describe how a system change with over time. They are widely used in fields such as physics, engineering, and mathematics. In many 
real-world problems, particularly those involving complex or nonlinear equations, it is often 
not possible to obtain exact solutions. That is why methods such as Euler’s method and 
Runge-Kutta are used to obtain approximate solutions. In this project, we examine a 
particular system consisting of three nonlinear equations known as the Lorenz-1960 system 
(i.e., one of the first meteorological models). The system is small, but complex enough to 
make a good testing ground for various solution methods. The Lorenz-1960 system is a 
great candidate for evaluating how well ANNs learn and predict behavior. Classical 
numerical methods address these problems step by step. Also they are typically accurate, but 
they are likely very slow and require more computational cost. We use numerical methods like Runge-Kutta to generate data that can be used to train artificial neural networks to learn how variables change over time.
The ANN learns to match its predictions with the results of the numerical solver during training. After training, the ANN can quickly predict results for any time without solving the equations again.
Therefore, the primary aim  of this project is to solve the Lorenz-1960 system 
and to see how well a plain ANN can learn its behavior and which ANN architecture provides
optimal results. 


\section{Motivation}
This project is motivated by both research and practical sides. Classic numerical methods are reliable and accurate. They work well for directly simulating equations. However, they can be slow when we need to run them many times. A trained artificial neural network (ANN), on the other hand, can solve the equation on its own. It makes predictions and acts like a copy of the real system. Many studies now use neural networks to solve differential equations, including physics-informed neural networks (PINNs), operator learning, and ANN-based approximators. This makes the topic both relevant and interesting to explore. Another reason for this project is the learning and hands-on benefits it offers. It helps build skills in solving ordinary differential equations (ODEs). It also teaches how to build and use neural network models. The project involves running experiments and studying the results. Comparing different methods is helpful. Even if simple ANNs are not as powerful as advanced methods, they help us understand their strengths and weaknesses.



\section{Objectives}
The primary objectives of this project are:

\begin{itemize}
    \item To solve the Lorenz equation, which is a coupled and nonlinear system.
\end{itemize}

\begin{itemize}
    \item To use ODE concepts, numerical methods like Runge-Kutta and SciPy solvers, and an ANN for solving the Lorenz equations.
\end{itemize}

\begin{itemize}
    \item To create accurate data using the Runge-Kutta method and use it to validate the ANN  performance and compare its results.
\end{itemize}

\begin{itemize}
    \item To compare different ANN architectures by changing the number of layers, the number of neurons per layer, and trying different activation functions to find the most accurate model.
\end{itemize}

\begin{itemize}
    \item To evaluate the model’s performance using mean squared error, learning stability, and computational cost, and see if a plain ANN can learn the Lorenz equation.
\end{itemize}

\begin{itemize}
    \item To compare the results with PINN based methods from other research to understand how this approach performs compared to recent techniques.
\end{itemize}


\section{Methodology}
We are following the “Waterfall Methodology.”Here, each phase is completed sequentially in a fixed order. It is suitable because our goals and requirements are clear from the start. Here, every step produces outputs that guide the next. This makes the process structured, simple to track, and easy to manage.

 First, during the Literature Review phase, we studied the Lorenz-1960 equation and its numerical solution methods, like RK4 and SciPy solvers. We also reviewed multiple ANN designs, activation functions, and training methods. 
 
  
  Next, in the Concept Development phase, the research problem and objectives were clearly defined. We focused on approximating the Lorenz-1960 equation using ANNs. We also worked on setting up the network layers, choosing the number of neurons, activation functions, learning rates, and training epochs, and generating training and validation data.
  
   
   In the Design phase, we developed the ANN models and training workflow, covering data preprocessing and optimized architecture. We established a numerical baseline using RK4 and SciPy to compare results. Assessment measures, visualization tools, and statistical tests were also defined to guarantee fair model comparison.
   
    
    In the Implementation phase, the numerical solvers and ANN models were implemented in Python . All models were trained on generated datasets . Accuracy, training time, number of parameters, and convergence were recorded. Each model was executed systematically according to the design.
    
    
    Then, in the Testing phase, we measured each ANN model's performance and predicted its performance using mean square error. We tested generalization on new data and analyzed the models' sensitivity to different parameter settings. The best models were compared with results from published studies.
   
    
    Finally, in the Deployment phase, we put all results, plots, and analyses into the final project report. We organized and documented the code to make it easy to reproduce our work. We included instructions for repeating the experiments and retraining the models. We also gave our final recommendations for the best ANN architecture. This phase ensures that the project is complete and ready.


For better understanding, a diagram of our methodology is given below:-

\begin{figure}[h]
    \centering
    \includegraphics[width=0.8\textwidth]{methodology.jpeg}
    \caption{Methodology Diagram}
    \label{fig:methodology}
\end{figure}


\section{Project Outcome}
The expected outcome of the project will be the clear and the justified 
determination of the most appropriate architecture for the ANN in estimating 
the solution to the Lorenz-1960 system of Ordinary Differential Equations 
(ODEs) within the scope of the experiment. This project must also include a 
numerical reference solution to the ODEs, the establishment of ANN 
baselines, and a comparative analysis of the results. \\[1pt]

The project will also contribute to the current understanding of the aptitude 
and limitations of ANN models in the approximation and solution of differential 
equations. If the ANN models are able to get a satisfactory level of accuracy in 
approximating the solution to the ODEs, the project will have demonstrated 
the capability of a simple ANN architecture to effectively act as an alternative 
numerical solver. This would suggest that a plain ANN, without the need for 
physics-informed loss terms or traditional mesh-dependent numerical 
methods, can on its own approximate ODE solutions with sufficient accuracy 
while being computationally less expensive in terms of inference time and the 
avoidance of iterative re-solving for each query. But if the models cannot 
achieve satisfactory accuracy in approximating the solution to the ODEs then 
the study will clearly outline the constraints of simple ANN architectures and 
underscore the need for advanced methods like PINNs. \\[1pt]

The expected outcome of the project will therefore be more than just the 
coding and experiment itself but will also involve a research based 
understanding of how ANN architectural choices like depth, width, and 
activation function can impact solution accuracy for coupled nonlinear ODE 
systems.


\section{Organization of the Report}
The rest of this report is outlined as follows. Chapter 2 outlines the theoretical 
background and literature review on Ordinary Differential Equations (ODE) 
solving, artificial neural networks, and existing neural methods for differential 
equations. Chapter 3 discusses about the project design, requirements, 
methodological plan, and task assignment. Chapter 4 outlines the 
implementation \& experiment design used in this research. Chapter 5 outlines 
the results and analysis on the tested ANN structures. Finally, Chapter 6 
concludes this report by focusing the main results, limitations, and potential 
directions for the future research. 
